\subsection*{3.7 Evaluating Statistical Claims: Observational Studies and Experiments }

\vspace{0.5em}

\setcounter{questionnum}{0}


\noindent\markforreview
\vspace{1em}

A sample of 40 fourth-grade students was selected at random from a certain school.  
The 40 students completed a survey about the morning announcements, and 32 thought the announcements were helpful.  
Which of the following is the largest population to which the results of the survey can be applied?

\vspace{0.75em}

\choicebox{A}{The 40 students who were surveyed}  
\choicebox{B}{All fourth-grade students at the school}  
\choicebox{C}{All students at the school}  
\choicebox{D}{All fourth-grade students in the county in which the school is located}

\vspace{2em}


\newpage
\noindent\markforreview
\vspace{0.5em}

A trivia tournament organizer wanted to study the relationship between the number of points a team scores in a trivia round  
and the number of hours that a team practices each week. For the study, the organizer selected 55 teams at random  
from all trivia teams in a certain tournament. The table displays the information for the 40 teams in the sample that practiced  
for at least 3 hours per week.


\begin{center}
    

\begin{tabular}{|l|c|c|c|}
\hline
Hours practiced & 6 to 13 points & 14 or more points & Total \\
\hline
3 to 5 hours & 6 & 4 & 10 \\
More than 5 hours & 4 & 26 & 30 \\
\hline
Total & 10 & 30 & 40 \\
\hline
\end{tabular}
\end{center}

\vspace{0.5em}

\noindent Which option is the largest population to which the results of the study can be generalized?



\choicebox{A}{All trivia teams in the tournament that scored 14 or more points in the round}  
\choicebox{B}{The 55 trivia teams in the sample}  
\choicebox{C}{The 40 trivia teams in the sample that practiced for at least 3 hours per week}  
\choicebox{D}{All trivia teams in the tournament}

\vspace{2em}

\noindent\markforreview
\vspace{1em}

A psychologist designed and conducted a study to determine whether playing a certain educational game  
increases middle school students’ accuracy in adding fractions.  
For the study, the psychologist chose a random sample of 35 students from all of the students  
at one of the middle schools in a large city. The psychologist found that students who played the game  
showed significant improvement in accuracy when adding fractions.  
What is the largest group to which the results of the study can be generalized?

\vspace{0.75em}

\choicebox{A}{The 35 students in the sample}  
\choicebox{B}{All students at the school}  
\choicebox{C}{All middle school students in the city}  
\choicebox{D}{All students in the city}

\vspace{2em}

\noindent\markforreview
\vspace{1em}

Near the end of a US cable news show, the host invited viewers to respond to a poll on the show’s website that asked,  
“Do you support the new federal policy discussed during the show?” At the end of the show, the host reported that  
28\% responded “Yes,” and 70\% responded “No.”  
Which of the following best explains why the results are unlikely to represent the sentiments of the population of the United States?

\vspace{0.75em}

\choicebox{A}{The percentages do not add up to 100\%, so any possible conclusions from the poll are invalid.}  
\choicebox{B}{Those who responded to the poll were not a random sample of the population of the United States.}  
\choicebox{C}{There were not 50\% “Yes” responses and 50\% “No” responses.}  
\choicebox{D}{The show did not allow viewers enough time to respond to the poll.}

\vspace{2em}

\noindent\markforreview
\vspace{1em}

To determine the mean number of children per household in a community,  
Tabitha surveyed 20 families at a playground.  
For the 20 families surveyed, the mean number of children per household was 2.4.  
Which of the following statements must be true?

\vspace{0.75em}

\choicebox{A}{The mean number of children per household in the community is 2.4.}  
\choicebox{B}{A determination about the mean number of children per household in the community should not be made because the sample size is too small.}  
\choicebox{C}{The sampling method is flawed and may produce a biased estimate of the mean number of children per household in the community.}  
\choicebox{D}{The sampling method is not flawed and is likely to produce an unbiased estimate of the mean number of children per household in the community.}






\end{document}